% Standalone LaTeX source. No external figures or bibliography files are needed.
% Adapted from Jamil Zhumabek, Matrix Sketching and Linear Programming (2025).
% Experiments are a newly executed, revised protocol; see PROVENANCE.md.
\documentclass[12pt,letterpaper,oneside]{report}
\usepackage[T1]{fontenc}
\usepackage[utf8]{inputenc}
\usepackage[margin=1in]{geometry}
\usepackage{amsmath,amssymb,amsthm}
\usepackage{array,booktabs}
\usepackage{capt-of}
\usepackage{graphicx}
\usepackage{enumitem}
\usepackage{needspace}
\usepackage{hyperref}
\hypersetup{hidelinks,pdftitle={Matrix Sketching and Linear Programming},pdfauthor={Jamil Zhumabek}}
\linespread{1.5}
\setlength{\parindent}{0.25in}
\setlength{\parskip}{0pt}
\setlength{\emergencystretch}{2em}
\clubpenalty=10000
\widowpenalty=10000
\newtheorem{theorem}{Theorem}[section]
\newtheorem{lemma}[theorem]{Lemma}
\theoremstyle{definition}
\newtheorem{definition}[theorem]{Definition}
\newcommand{\R}{\mathbb{R}}
\newcommand{\E}{\mathbb{E}}
\newcommand{\norm}[1]{\left\lVert #1\right\rVert}
\newcommand{\feas}{\operatorname{feas}}
\newcommand{\negmass}{\operatorname{neg}}
\newcommand{\obj}{\operatorname{obj}}
\begin{document}

\begin{titlepage}
\centering
\vspace*{0.6in}
{\Large\bfseries Matrix Sketching and Linear Programming\par}
\vspace{0.45in}
by\par
{\large Jamil Zhumabek\par}
\vspace{0.8in}
Algorithms for Big Data\par
Course Project\par
\vspace{0.6in}
Mohamed bin Zayed University of Artificial Intelligence\par
\vspace{0.4in}
October 2026\par
\vfill
\begin{minipage}{0.85\textwidth}
\small
This report is adapted from the author's M.Sc. thesis, \emph{Matrix
Sketching and Linear Programming}, Nazarbayev University, 2025. It
retains the part on random projections for linear programming and presents
a fresh, reduced experimental evaluation with a Netlib benchmark extension.
\end{minipage}
\end{titlepage}
\setcounter{page}{2}

\chapter*{Abstract}
Random projection has become a standard tool for dimensionality reduction,
yet research is still being conducted on its efficiency for large linear
programs. This project studies the random-projection approach to linear
programming, develops the relevant theoretical background, and evaluates
the method through numerical experiments.

After reviewing the Johnson--Lindenstrauss lemma, we prove that if the
original linear programming problem is unbounded, then the projected one
is unbounded as well. A counterexample shows that the converse is not
necessarily true. The forward implication holds for every linear
projection because the projected constraints relax the original ones.

We conduct numerical experiments on synthetically generated feasible and
infeasible instances, and on five established Netlib LP benchmarks. The experiments compare running time, equality
residuals, nonnegativity violations, and objective discrepancies. They
distinguish a solution of the projected problem from a feasible solution
of the original problem. In particular, reducing the number of constraints
does not by itself guarantee either a computational improvement or a
valid recovered solution. The Netlib extension also measures matrix storage
and checks unbounded-ray certificates. It evaluates the basic constraint
relaxation, not sketching methods for optimization in general.

\chapter{Introduction}
\section{Motivation and Background}
Matrix sketching is a technique designed to efficiently approximate a
large matrix using a more compact representation. The goal is to preserve
properties of interest within a guaranteed approximation ratio. Random
projections achieve this by randomly combining rows of the original
matrix. The Johnson--Lindenstrauss lemma provides a fundamental guarantee
for preserving distances between a finite set of points in a lower
dimensional space.

In this work we concentrate particularly on the random projections
method with application to linear programming problems. Linear
programming is widely applied in operations research, machine learning,
economics, and combinatorics. However, a large constraint matrix can make
its solution computationally expensive. A smaller projected matrix
provides a possible way of reducing this cost.

Our research problem is to understand what is retained when the equality
constraints of a linear program are projected, and whether the reduction
in problem size improves computational performance. We follow the
random-projection formulation studied by Vu, Poirion, and Liberti
\cite{vu2018}. This report focuses on its basic projected relaxation and
the two recovery procedures described in the author's thesis
\cite{thesis2025}. It does not
claim the full approximation guarantees of \cite{vu2018}, which require
additional assumptions and restrictions on the projected problem.

\section{Relationship to Previous Work}
The theoretical starting point is the author's M.Sc. thesis,
\emph{Matrix Sketching and Linear Programming} \cite{thesis2025}.
The earlier experimental implementation is available in the companion
GitHub project \cite{legacyrepo}. This course project retains the
relevant theoretical material and the two recovery strategies, while
replacing the earlier experimental chapter with newly executed
experiments. The unboundedness result is retained from the thesis;
no claim of a new literature contribution is made here.

The previous repository is cited for provenance and is not modified.
The revised code, recorded seeds, raw observations, environment metadata,
and report source are supplied in the separate course-project repository
\href{https://github.com/Jamil997/random-projections-lp-course-project}{\texttt{random-projections-lp-course-project}}. The master's thesis PDF
and previous notebook outputs are not part of this package.

\section{Research Objectives}
The primary objectives of this project are as follows:
\begin{enumerate}[label=\arabic*)]
\item To introduce random projections and the theoretical background
needed for their application to linear programming.
\item To prove the preservation of unboundedness under projection and
show that the converse does not necessarily hold.
\item To conduct numerical experiments on feasible and infeasible
instances and established LP benchmarks, comparing solution quality,
running time, and projected matrix storage.
\end{enumerate}

\section{Matrix Sketching and Random Projections}
Let $A$ be an $m\times n$ constraint matrix. We consider the following
linear programming problem:
\begin{equation}\label{eq:original}
\begin{aligned}
(P)\qquad &\text{minimize} &&c^\top x\\
&\text{subject to} &&Ax=b,\\
&&&x\geq0,
\end{aligned}
\end{equation}
where $c\in\R^n$ and $b\in\R^m$.
Let $S$ be a short and fat $k\times m$ matrix, where $k<m$. For a
Gaussian random projection we take independent entries
$S_{ij}\sim\mathcal{N}(0,1/k)$. The matrix product has dimensions
\[
\underbrace{S}_{k\times m}\;
\underbrace{A}_{m\times n}
=\underbrace{SA}_{k\times n}.
\]
Then, we solve the projected problem:
\begin{equation}\label{eq:projected}
\begin{aligned}
(P_S)\qquad &\text{minimize} &&c^\top x\\
&\text{subject to} &&SAx=Sb,\\
&&&x\geq0.
\end{aligned}
\end{equation}
The number of equality constraints decreases from $m$ to $k$, while the
number of variables remains $n$. We compare solution quality and running
time between the original and the sketched problem.

\section{Other Sketching Methods and Their Relevance}
The thesis surveys several sketching techniques before specializing to
linear programming \cite[pp.~6--16]{thesis2025}. We retain a focused part
of that background here. For a row sketch $B=SA$, one possible quality
criterion is the covariance error $\norm{A^\top A-B^\top B}_2$.
For a unit vector $z$, this controls the error in $\norm{Az}_2^2$.
Such a criterion concerns matrix geometry; by itself it does not certify
the feasibility or boundedness of a projected LP.

\subsection{Row Sampling and Sparse Embeddings}
Let $a_i^\top$ be row $i$ of $A$, with $A\ne0$. Norm-based row sampling
draws $k$ independent indices with probabilities
$p_i=\norm{a_i}_2^2/\norm{A}_F^2$ and takes the sampled rows
$a_i^\top/\sqrt{kp_i}$. Zero rows are never sampled. The normalization gives
\[
\E[B^\top B]=\sum_{j=1}^k\sum_{i:p_i>0}
p_i\frac{a_i a_i^\top}{kp_i}=A^\top A.
\]
Unbiasedness is not a uniform approximation guarantee, and leverage-based
sampling is a different choice of sampling distribution.

CountSketch instead hashes each original row to a bucket
$h(i)\in\{1,\ldots,k\}$ and assigns an independent random sign
$s_i\in\{-1,1\}$. Its sketching matrix has exactly one nonzero per column:
\[
S_{ri}=s_i\mathbf{1}\{h(i)=r\},\qquad
B_{r,:}=\sum_{i:h(i)=r}s_i A_{i,:}.
\]
For a fixed vector $v$, independent mean-zero signs give
$\E\norm{Sv}_2^2=\norm{v}_2^2$. Stronger embedding guarantees require
dimension and probability conditions \cite{clarkson2013}; this expectation
alone is not the JL lemma. Accumulating the contributions can take time
proportional to $\operatorname{nnz}(A)$, apart from allocation and storage
management. Collisions may cancel entries or leave empty sketch rows.
The Netlib experiments test CountSketch alongside Gaussian and
sparse-valued random projections, using actual sparse storage.

\subsection{Low-Rank Approximation and Streaming Sketches}
Randomized range finding draws $\Omega\in\R^{n\times\ell}$, computes
$Y=A\Omega$, and obtains an orthonormal basis $Q$ for its range.
A factorization of the smaller matrix $Q^\top A$ then approximates a
factorization of $A$ \cite{halko2011}. This data-dependent low-rank
approximation is distinct from replacing $Ax=b$ by an oblivious row sketch.
Frequent Directions is a deterministic streaming approach that shrinks
singular values while maintaining a small covariance sketch
\cite{ghashami2016}. These methods provide context from the thesis;
neither is implemented or experimentally evaluated in this project.

\subsection{What Would Preserve All Equality Residuals?}
Suppose that $S$ is a subspace embedding for
$\mathcal{V}=\operatorname{range}([A\ b])$ with $0<\epsilon<1$, so that
\[
(1-\epsilon)\norm{v}_2^2\leq\norm{Sv}_2^2
\leq(1+\epsilon)\norm{v}_2^2\qquad(v\in\mathcal{V}).
\]
Every residual $Ax-b$ belongs to $\mathcal{V}$. The lower bound implies
$S(Ax-b)=0$ if and only if $Ax-b=0$. Thus this stronger condition would
preserve all equalities exactly. However, it requires $S$ to be injective
on $\mathcal{V}$ and hence $k\geq\dim\mathcal{V}$. If $A$ has full row
rank, $\mathcal{V}=\R^m$ and no sketch with $k<m$ can satisfy it.
This elementary observation explains why preserving a finite set of
distances cannot automatically justify preservation of an entire LP
feasible set. It is not a claimed new theorem.

\subsection{Sketching as a Preconditioner}
The thesis also reviews randomized interior-point methods
\cite[pp.~31--35]{thesis2025}. Their computational difficulty includes
repeated linear systems involving a normal matrix
$M=AD^2A^\top$, where $D$ is a positive diagonal scaling.
A suitable sketch $W$ can form
$\widetilde M=ADWW^\top DA^\top$ as a preconditioner for an iterative
solve. In this use, sketching accelerates a linear-algebra subproblem;
it does not simply discard equality constraints from the optimization
problem. The analysis of Chowdhury et al. includes safeguards for the
inexact solves \cite{chowdhury2022}. We do not implement that algorithm.
Consequently, a failure of our basic relaxation is not evidence against
randomized preconditioning or all sketch-based LP methods.

\chapter{Main Result}
This chapter establishes the theoretical background for analyzing linear
programming problems under random projections. The
Johnson--Lindenstrauss lemma explains the geometric motivation for
dimensionality reduction. Farkas' lemma gives a certificate of
infeasibility. We then prove preservation of unboundedness, retaining the
result and counterexample developed in the thesis \cite{thesis2025}.

\section{Preliminaries}
\begin{lemma}[The Johnson--Lindenstrauss lemma]\label{lem:jl}
Let $V\subset\R^m$ be a set of $N\geq2$ points and let $0<\epsilon<1$.
For every integer
\begin{equation}\label{eq:jl-dimension}
k\geq\frac{4\ln N}{\epsilon^2/2-\epsilon^3/3},
\end{equation}
there exists a linear map $f:\R^m\to\R^k$ such that, for all $u,v\in V$,
\[
(1-\epsilon)\norm{u-v}_2^2
\leq\norm{f(u)-f(v)}_2^2
\leq(1+\epsilon)\norm{u-v}_2^2.
\]
\end{lemma}

\begin{proof}
Our aim is to estimate the length of a fixed vector after it has been
mapped into a lower dimensional space. We use the Gaussian concentration
argument underlying elementary proofs of the lemma
\cite{dasgupta2003,achlioptas2003}.

Let the entries of $S\in\R^{k\times m}$ be independent
$\mathcal{N}(0,1/k)$ variables, and define $f(w)=Sw$. For a fixed
$w\ne0$, the random variable
\[
Z=\frac{k\norm{Sw}_2^2}{\norm{w}_2^2}
\]
is the sum of squares of $k$ independent standard normal variables. Its
moment generating function is
\[
\E[e^{tZ}]=(1-2t)^{-k/2},\qquad t<\tfrac12.
\]
Applying Markov's inequality to $e^{tZ}$ and minimizing the resulting
bound over $t>0$, with $t=\epsilon/[2(1+\epsilon)]$, gives
\begin{align*}
\Pr[Z\geq k(1+\epsilon)]
&\leq\exp\!\left[-\frac{k}{2}
\bigl(\epsilon-\ln(1+\epsilon)\bigr)\right]\\
&\leq\exp\!\left[-\frac{k}{2}
\left(\frac{\epsilon^2}{2}-\frac{\epsilon^3}{3}\right)\right].
\end{align*}
The last step follows from
$\ln(1+\epsilon)\leq\epsilon-\epsilon^2/2+\epsilon^3/3$.
Similarly, applying Markov's inequality to $e^{-tZ}$ and choosing
$t=\epsilon/[2(1-\epsilon)]$ gives
\begin{align*}
\Pr[Z\leq k(1-\epsilon)]
&\leq\exp\!\left[\frac{k}{2}
\bigl(\epsilon+\ln(1-\epsilon)\bigr)\right]\\
&\leq\exp(-k\epsilon^2/4).
\end{align*}
Under \eqref{eq:jl-dimension}, each tail probability is at most $N^{-2}$.
Apply these estimates to each nonzero difference $u-v$, and then use the
union bound over the $\binom{N}{2}$ pairs. The probability that any pair
violates the stated inequalities is at most
\[
\binom{N}{2}\frac{2}{N^2}=1-\frac1N<1.
\]
Hence, there exists a mapping with the required properties.
\end{proof}

The lemma guarantees distance preservation for a finite set of vectors.
It does not automatically guarantee that all feasible points, all
recession directions, or the optimal value of an LP are preserved.
The analysis of \cite{vu2018} accounts for additional geometric
conditions. Our experiments use $k=\lfloor m/2\rfloor$ as a fixed
compression level; this choice is empirical and is not asserted to satisfy
\eqref{eq:jl-dimension} for a prescribed $\epsilon$.

\begin{lemma}[Farkas' lemma]\label{lem:farkas}
Exactly one of the following two systems has a solution:
\[
Ax=b,\quad x\geq0;
\qquad\text{or}\qquad
A^\top y\geq0,\quad b^\top y<0.
\]
\end{lemma}
This is the standard equality-form alternative for linear programming
\cite{vanderbei2020}. The two systems cannot both hold: if they did, then
\[
b^\top y=(Ax)^\top y=x^\top A^\top y\geq0,
\]
contradicting $b^\top y<0$. We use this certificate to generate infeasible
instances in Chapter~\ref{chap:experiments}.

\begin{definition}[Recession direction]
For a nonempty standard-form polyhedron
$F=\{x\in\R^n:Ax=b,\ x\geq0\}$, a recession direction is a vector
$d$ satisfying $Ad=0$ and $d\geq0$. Then $x+td\in F$ for every
$x\in F$ and $t\geq0$.
\end{definition}

\section{Preservation of Unboundedness}
\begin{theorem}\label{thm:unbounded}
Assume the original LP
\[
\min\{c^\top x:Ax=b,\ x\geq0\}
\]
is unbounded below. Then for every matrix $S\in\R^{k\times m}$ the
projected LP
\[
\min\{c^\top x:SAx=Sb,\ x\geq0\}
\]
is unbounded below as well. The converse is not necessarily true.
\end{theorem}

\begin{proof}
Since the original LP is unbounded, there exists a feasible point $x_0$
and a recession direction $d$ such that
\[
Ax_0=b,\quad x_0\geq0,\qquad
Ad=0,\quad d\geq0,\quad c^\top d<0.
\]
This is the standard recession characterization of unboundedness for a
nonempty linear program \cite{vanderbei2020}.
To construct an unbounded ray, define for any $t\geq0$
\[
x(t)=x_0+td.
\]
Since $Ad=0$, substituting into the projected constraints gives
\[
SA(x_0+td)=SAx_0+tSAd=Sb.
\]
Also, $x_0+td\geq0$. In other words, the same ray is feasible for the
projected problem. The objective remains unchanged:
\[
c^\top(x_0+td)=c^\top x_0+t c^\top d,
\]
which tends to $-\infty$ as $t\to\infty$. Hence, the projected LP is
unbounded. No condition on the distribution of $S$ is required because
$Ad=0$ implies $SAd=0$ identically.

Let us show that the converse does not hold using the following
counterexample. For the original LP, let
\[
A=\begin{pmatrix}1&0\\0&1\end{pmatrix},\qquad
b=\begin{pmatrix}1\\2\end{pmatrix},\qquad
c=\begin{pmatrix}0\\-1\end{pmatrix}.
\]
The solution of $Ax=b$ is $x=(1,2)^\top$. Thus, the feasible set is a
single point and the original optimal value is $-2$.
For the projected LP, let
\[
S=\begin{pmatrix}1&0\end{pmatrix}.
\]
Then $SA=(1\ \ 0)$ and $Sb=1$. The only equality is $x_1=1$, while
$x_2\geq0$ is free to increase. The objective $c^\top x=-x_2$ tends to
$-\infty$ as $x_2\to\infty$. Thus, the projected LP is unbounded, yet
the original LP is bounded.
\end{proof}

The theorem can also be seen directly from inclusion of feasible sets:
projection relaxes the equalities while keeping the objective and
nonnegativity constraints. In particular, $v(P_S)\leq v(P)$ whenever
both optimal values are finite. This elementary observation is useful for
interpreting the experiments; it is not a probabilistic consequence of
the Johnson--Lindenstrauss lemma.

\chapter{Synthetic Experiments and Results}\label{chap:experiments}
\section{Experimental Setup}
We considered the pairs $(m,n)=(100,150),(250,375),(500,750)$,
constraint matrix densities $\mathrm{dens}\in\{0.1,0.5\}$, and
$k=\lfloor m/2\rfloor$. For each configuration and each of the three
instance families, we generated 10 independent original LPs. Entries
of $A$ were sampled uniformly on $[0,1]$ and retained independently with
probability $\mathrm{dens}$; other entries were set to zero.
The objective vector was always $c=\mathbf{1}$.

\begin{enumerate}[label=\arabic*)]
\item \emph{Feasible instances:} sample $x_0$ uniformly on $[0,1]^n$
and set $b=Ax_0$.
\item \emph{Legacy all-negative instances:} sample
$b_i=-u_i$, independently, with $u_i\sim U[0,1]$.
This retains the right-hand-side distribution in the previous
infeasible-instance notebook \cite{legacyrepo}.
\item \emph{Additional mixed-sign stress test:} sample each $b_i$
uniformly on $[0.5,1.5]$ and replace $b_1$ by $-b_1$.
This is an additional instance family, not the original thesis generator.
\end{enumerate}
For both infeasible families, $A^\top e_1\geq0$ and
$b^\top e_1<0$ in every generated instance. Thus, Lemma~\ref{lem:farkas}
certifies infeasibility before solving. Both families contain an immediate
contradiction in the original equalities, which presolve may detect.

This is a fresh execution of an explicitly revised protocol. The previous
notebooks \cite{legacyrepo} were inspected at commit
\texttt{3d2ff96}; the full commit identifier is recorded in the
repository's provenance file. Their saved configurations used larger
matrices, different compression, and fixed-count nonzero selection.
The updated experiment uses the smaller grid above and Bernoulli masks.
For feasible instances, it uses the same explicit solver for the original
and projected LPs, replacing the earlier combination of SciPy and
CVXPY's automatic solver choice. Direct and SVD solves replace explicit
matrix inversion and normal equations. Projector normalizations,
random streams, and solver-status reporting are made explicit below.
Therefore, these are new measurements of the stated protocol, not an
exact reconstruction of the old timing tables. No earlier experimental
measurements are reused in this chapter.

\subsection{Random Projectors}
For feasible instances, we used a Gaussian projector with independent
entries $\mathcal{N}(0,1/k)$. For infeasible instances, we used the
following two projectors:
\begin{enumerate}[label=\arabic*)]
\item A sparse random projector, with independent entries
\[
S_{ij}=\sqrt{\frac3k}
\begin{cases}
+1,&\text{with probability }1/6,\\
0,&\text{with probability }2/3,\\
-1,&\text{with probability }1/6.
\end{cases}
\]
This normalization gives $\E[S_{ij}^2]=1/k$ \cite{achlioptas2003}.
\item An orthogonal random projector. We formed a Gaussian
$m\times k$ matrix and computed a reduced QR decomposition, obtaining
$Q^\top Q=I_k$. We set $S=\sqrt{m/k}\,Q^\top$.
\end{enumerate}
The scaling factors are needed for the usual distance-preservation
interpretation. Multiplying all of $S$ by the same nonzero scalar leaves
the equality system $SAx=Sb$ unchanged in exact arithmetic.

\subsection{Recovery of a Primal Vector}
We retained the following two recovery procedures from the thesis
\cite{thesis2025}, using the original implementation \cite{legacyrepo}
as a reference for the method.
Both return a candidate vector that must be checked against the original
constraints.

\noindent\textbf{Algorithm 1: Exact-solve (Dual-based recovery)}
\begin{enumerate}[label=\arabic*:,nosep]
\item Obtain a projected dual solution $\lambda^*\in\R^k$ and form
\[
y_{\mathrm{prox}}=S^\top\lambda^*,\qquad
z_j=\frac{c_j-A_j^\top y_{\mathrm{prox}}}{\norm{A_j}_2}.
\]
Here $\lambda^*$ is the SciPy equality marginal, whose sign convention
gives $c-(SA)^\top\lambda^*\geq0$.
Zero column norms are assigned an infinite score, and ties are resolved
by column index.
\item Select $B\subseteq\{1,\ldots,n\}$ as the indices of the $m$
smallest normalized reduced costs $z_j$.
\item If $A_B$ is nonsingular, solve $A_Bx_B=b$ and set $x_j=0$ for
$j\notin B$. A singular selected matrix is recorded as a recovery failure.
\end{enumerate}

\Needspace{5\baselineskip}
\noindent\textbf{Algorithm 2: Pseudoinverse-solve (Primal-based recovery)}
\begin{enumerate}[label=\arabic*:,nosep]
\item Obtain a projected primal solution $\bar x\in\R^n$.
\item Select $H$ as the indices of the $m$ largest entries of $\bar x$.
\item Compute $x_H=A_H^+b$ and set $x_j=0$ for $j\notin H$.
\end{enumerate}
Here $A_H^+$ denotes the Moore--Penrose pseudoinverse. We use an SVD
least-squares solve instead of explicitly forming
$(A_H^\top A_H)^+A_H^\top$, which is algebraically equivalent but can
worsen numerical conditioning. Neither procedure enforces $x\geq0$.

\subsection{Solution Quality and Running Time}
For a returned vector $x$, we use the thesis's normalized metrics:
\begin{align*}
\feas(x)&=\frac{\norm{Ax-b}_1}{\norm{b}_1},\\
\negmass(x)&=\frac{\sum_{j:x_j<0}|x_j|}
{\norm{x}_1},\\
\obj(x)&=\frac{|c^\top x-v(P)|}{|v(P)|}.
\end{align*}
We evaluate these metrics separately for the projected solution $\bar x$
and the two recovered vectors. For $\bar x$, the objective is the
projected optimal value. For a recovered vector, $c^\top x$ is only a
candidate objective value unless original feasibility has been checked.
A small equality residual does not compensate for negative components.
The generated right-hand sides and the feasible reference optimal values
are nonzero. The implementation leaves a metric undefined if its
reference denominator vanishes, except that $\negmass(0)=0$.

A candidate is accepted as feasible for the original LP only if
\[
\norm{Ax-b}_\infty\leq 10^{-7}\max(1,\norm b_\infty),
\qquad \min_j x_j\geq-10^{-7}.
\]
Timing tables give means and, for the feasible family, sample standard
deviations over 10 instances. Quality tables give arithmetic means;
counts refer to the same 10 instances. All sample standard deviations,
solver messages, timing components, and seeds are supplied in the raw
and summary files.

The experiments were executed on 13 September 2026 using Python 3.12.14,
NumPy 2.5.3, SciPy 1.18.1, and HiGHS 1.12.0 on an ARM64 Mac with
Apple Accelerate. We used HiGHS dual simplex (\texttt{highs-ds}),
presolve, and primal and dual feasibility tolerances of $10^{-7}$.
Single-thread settings were requested for HiGHS and numerical libraries;
Accelerate's active thread count was not independently measured.
Separate data and projector streams were derived from master seed
20260913. The two infeasible projectors share each original instance
within each family. In total, there are 180 independent original LPs
and 300 projected LPs.

We measured elapsed wall-clock time. Original time covers its solver
call. Projected time includes sampling $S$, computing $SA$ and $Sb$,
and solving the projected LP. Each recovery total includes that common
projected cost and its own recovery step. Data generation, metric
evaluation, and diagnostic retries are excluded. All matrices, including
the sparse-valued projector, were stored densely. This is therefore not
a sparse-storage benchmark.

\section{Feasible Instances}
All original and projected feasible-family LPs were solved to optimality.
Table~\ref{tab:fresh-feasible-times} reports the time to solve the original problem,
the projected relaxation, and each projection-and-recovery workflow.
A shorter projected workflow must be interpreted together with the
quality checks.

% BEGIN GENERATED feasible_times
\par\medskip
\noindent\begin{minipage}{\linewidth}
\small\linespread{1}\selectfont\centering
\setlength{\tabcolsep}{4pt}
\renewcommand{\arraystretch}{1.08}
\captionof{table}{Fresh feasible-instance timings in milliseconds (mean $\pm$ sample standard deviation; $d$: density). Projected totals include sampling, multiplication and solving; recovery totals add the indicated recovery.}\label{tab:fresh-feasible-times}
\medskip
\sbox0{%
\begin{tabular}{rr rrrr}
\toprule
$m$ & $d$ & Original & Projected & Dual total & Primal total \\
\midrule
100 & 0.1 & $3.05\pm 0.29$ & $2.86\pm 0.13$ & $3.04\pm 0.19$ & $3.56\pm 0.13$ \\
100 & 0.5 & $6.39\pm 0.22$ & $2.85\pm 0.15$ & $3.00\pm 0.15$ & $3.55\pm 0.14$ \\
250 & 0.1 & $40.95\pm 1.64$ & $25.40\pm 0.74$ & $25.97\pm 0.77$ & $28.65\pm 0.76$ \\
250 & 0.5 & $85.65\pm 4.31$ & $26.99\pm 1.31$ & $27.54\pm 1.31$ & $30.26\pm 1.31$ \\
500 & 0.1 & $481.54\pm 37.79$ & $206.13\pm 5.61$ & $208.28\pm 5.61$ & $217.63\pm 5.63$ \\
500 & 0.5 & $886.02\pm 65.74$ & $217.90\pm 9.64$ & $220.07\pm 9.67$ & $229.43\pm 9.62$ \\
\bottomrule
\end{tabular}
}%
\ifdim\wd0>\linewidth
  \resizebox{\linewidth}{!}{\usebox0}
\else
  \usebox0
\fi
\end{minipage}
\par\medskip
% END GENERATED feasible_times

Table~\ref{tab:fresh-feasible-quality} reports original-system residuals, negative-mass
ratios, and objective discrepancies. The projected solutions retain
nonnegativity but do not satisfy the original equalities. The recovered
vectors have very small equality residuals, but substantial negative
components. Neither recovery procedure returns an accepted feasible
solution in this benchmark. Accordingly, a recovery timing is not the
time to obtain a valid solution of the original LP.

% BEGIN GENERATED feasible_quality
\par\medskip
\noindent\begin{minipage}{\linewidth}
\small\linespread{1}\selectfont\centering
\setlength{\tabcolsep}{4pt}
\renewcommand{\arraystretch}{1.08}
\captionof{table}{Mean relative equality residual, negative mass ratio and relative objective gap ($d$: density; Raw: projected primal; Dual/Primal: Algorithms 1/2). Accepted counts require equality feasibility and nonnegativity.}\label{tab:fresh-feasible-quality}
\medskip
\sbox0{%
\begin{tabular}{rr l rrr r}
\toprule
$m$ & $d$ & Method & Eq. residual & Neg. mass & Obj. gap & Accepted \\
\midrule
100 & 0.1 & Raw & $2.98\times 10^{-1}$ & 0.000 & 0.209 & 0/10 \\
100 & 0.1 & Dual & $3.52\times 10^{-14}$ & 0.444 & 0.799 & 0/10 \\
100 & 0.1 & Primal & $2.22\times 10^{-15}$ & 0.417 & 0.236 & 0/10 \\
100 & 0.5 & Raw & $1.36\times 10^{-1}$ & 0.000 & 0.106 & 0/10 \\
100 & 0.5 & Dual & $1.52\times 10^{-15}$ & 0.446 & 0.066 & 0/10 \\
100 & 0.5 & Primal & $4.94\times 10^{-14}$ & 0.388 & 1.717 & 0/10 \\
250 & 0.1 & Raw & $2.07\times 10^{-1}$ & 0.000 & 0.163 & 0/10 \\
250 & 0.1 & Dual & $2.76\times 10^{-14}$ & 0.454 & 0.837 & 0/10 \\
250 & 0.1 & Primal & $5.25\times 10^{-15}$ & 0.432 & 0.303 & 0/10 \\
250 & 0.5 & Raw & $8.40\times 10^{-2}$ & 0.000 & 0.070 & 0/10 \\
250 & 0.5 & Dual & $4.94\times 10^{-15}$ & 0.451 & 0.083 & 0/10 \\
250 & 0.5 & Primal & $2.57\times 10^{-15}$ & 0.444 & 0.085 & 0/10 \\
500 & 0.1 & Raw & $1.59\times 10^{-1}$ & 0.000 & 0.126 & 0/10 \\
500 & 0.1 & Dual & $8.80\times 10^{-14}$ & 0.474 & 0.557 & 0/10 \\
500 & 0.1 & Primal & $9.09\times 10^{-15}$ & 0.458 & 0.280 & 0/10 \\
500 & 0.5 & Raw & $6.50\times 10^{-2}$ & 0.000 & 0.052 & 0/10 \\
500 & 0.5 & Dual & $1.56\times 10^{-14}$ & 0.465 & 0.074 & 0/10 \\
500 & 0.5 & Primal & $2.11\times 10^{-15}$ & 0.456 & 0.046 & 0/10 \\
\bottomrule
\end{tabular}
}%
\ifdim\wd0>\linewidth
  \resizebox{\linewidth}{!}{\usebox0}
\else
  \usebox0
\fi
\end{minipage}
\par\medskip
% END GENERATED feasible_quality

These observations are consistent with the recovery procedures:
solving the selected linear system does not enforce nonnegativity.
Moreover, a basic projected optimum has at most $k$ positive entries.
Algorithm 2 selects $m>k$ entries, so remaining columns may be selected
among zero-valued ties. We resolve such ties by ascending column index,
which can influence the recovered vector. This experiment identifies a
limitation of these heuristics for the selected instance distributions
and compression level, rather than a failure of every recovery method
for random-projection LPs.

\section{Infeasible Instances}
Every original LP in both infeasible families was reported infeasible.
We report primary projected statuses separately: I means infeasible;
F means optimal and hence feasible for the projected LP; E means an
unresolved solver error. A projected feasible status must not be
interpreted as feasibility of the original LP.

\subsection{Legacy All-Negative Family}
Table~\ref{tab:fresh-legacy-status} evaluates the original notebook's all-negative
right-hand-side distribution on the updated grid. All projected LPs
in this family were reported infeasible. This is an empirical observation
for the sampled instances, not a general preservation guarantee.

% BEGIN GENERATED legacy_status
\par\medskip
\noindent\begin{minipage}{\linewidth}
\small\linespread{1}\selectfont\centering
\setlength{\tabcolsep}{4pt}
\renewcommand{\arraystretch}{1.08}
\captionof{table}{Fresh all-negative-family results (legacy right-hand-side distribution). Here $d$ denotes density; times are means in milliseconds. Counts are primary projected-solver outcomes; all original LPs are infeasible.}\label{tab:fresh-legacy-status}
\medskip
\sbox0{%
\begin{tabular}{rr l rr rrr}
\toprule
$m$ & $d$ & Projector & Orig. (ms) & Proj. (ms) & Inf. & Opt. & Error \\
\midrule
100 & 0.1 & Sparse & 0.64 & 1.95 & 10 & 0 & 0 \\
100 & 0.1 & Orthogonal & 0.64 & 2.08 & 10 & 0 & 0 \\
100 & 0.5 & Sparse & 1.04 & 1.71 & 10 & 0 & 0 \\
100 & 0.5 & Orthogonal & 1.04 & 1.95 & 10 & 0 & 0 \\
250 & 0.1 & Sparse & 1.55 & 12.46 & 10 & 0 & 0 \\
250 & 0.1 & Orthogonal & 1.55 & 13.35 & 10 & 0 & 0 \\
250 & 0.5 & Sparse & 4.61 & 10.97 & 10 & 0 & 0 \\
250 & 0.5 & Orthogonal & 4.61 & 12.02 & 10 & 0 & 0 \\
500 & 0.1 & Sparse & 4.51 & 67.39 & 10 & 0 & 0 \\
500 & 0.1 & Orthogonal & 4.51 & 69.31 & 10 & 0 & 0 \\
500 & 0.5 & Sparse & 18.63 & 65.71 & 10 & 0 & 0 \\
500 & 0.5 & Orthogonal & 18.63 & 64.72 & 10 & 0 & 0 \\
\bottomrule
\end{tabular}
}%
\ifdim\wd0>\linewidth
  \resizebox{\linewidth}{!}{\usebox0}
\else
  \usebox0
\fi
\end{minipage}
\par\medskip
% END GENERATED legacy_status

\subsection{Additional Mixed-Sign Stress Test}
Table~\ref{tab:fresh-stress-status} shows a different outcome for the additional
mixed-sign family. The sparse projector returned 50 infeasible,
8 optimal, and 2 unresolved outcomes; the orthogonal projector returned
49 infeasible, 9 optimal, and 2 unresolved outcomes. Thus, some projected
relaxations are feasible even though their original LPs have an explicit
Farkas certificate. The contrast with the all-negative family shows why
the infeasible-instance generator needs to be specified.

% BEGIN GENERATED stress_status
\par\medskip
\noindent\begin{minipage}{\linewidth}
\small\linespread{1}\selectfont\centering
\setlength{\tabcolsep}{4pt}
\renewcommand{\arraystretch}{1.08}
\captionof{table}{Fresh mixed-sign stress-test results (additional family). Here $d$ denotes density; times are means in milliseconds. Counts are primary projected-solver outcomes; all original LPs are infeasible.}\label{tab:fresh-stress-status}
\medskip
\sbox0{%
\begin{tabular}{rr l rr rrr}
\toprule
$m$ & $d$ & Projector & Orig. (ms) & Proj. (ms) & Inf. & Opt. & Error \\
\midrule
100 & 0.1 & Sparse & 0.68 & 2.96 & 3 & 7 & 0 \\
100 & 0.1 & Orthogonal & 0.68 & 3.06 & 3 & 7 & 0 \\
100 & 0.5 & Sparse & 1.06 & 2.38 & 10 & 0 & 0 \\
100 & 0.5 & Orthogonal & 1.06 & 2.54 & 10 & 0 & 0 \\
250 & 0.1 & Sparse & 1.55 & 29.61 & 9 & 1 & 0 \\
250 & 0.1 & Orthogonal & 1.55 & 28.97 & 8 & 2 & 0 \\
250 & 0.5 & Sparse & 4.65 & 16.25 & 10 & 0 & 0 \\
250 & 0.5 & Orthogonal & 4.65 & 18.04 & 10 & 0 & 0 \\
500 & 0.1 & Sparse & 4.51 & 178.10 & 9 & 0 & 1 \\
500 & 0.1 & Orthogonal & 4.51 & 192.98 & 10 & 0 & 0 \\
500 & 0.5 & Sparse & 18.90 & 86.32 & 9 & 0 & 1 \\
500 & 0.5 & Orthogonal & 18.90 & 102.85 & 8 & 0 & 2 \\
\bottomrule
\end{tabular}
}%
\ifdim\wd0>\linewidth
  \resizebox{\linewidth}{!}{\usebox0}
\else
  \usebox0
\fi
\end{minipage}
\par\medskip
% END GENERATED stress_status

The four unresolved primary runs returned an unknown HiGHS model status.
Separate diagnostic solves with HiGHS interior point classified all four
as infeasible; disabling presolve for dual simplex did not resolve them.
These checks are saved separately and do not replace the primary counts
or enter the benchmark timings.

Across both infeasible families, the original LP was faster on average
than constructing and solving its projected relaxation in every
configuration. Projection adds sampling and matrix-product costs and
usually makes the constraint matrix dense. It can also hide an immediate
presolve contradiction, or produce a feasible relaxation. These results
concern intentionally easy-to-certify infeasible instances and are not
evidence that projection is slower for arbitrary infeasible LPs.

\section{Verification and Reproducibility}
Since $c=\mathbf{1}$ and $x\geq0$, every feasible LP in the random
benchmark is bounded below by zero. These runs do not test
unboundedness. We therefore performed separate deterministic checks.

The counterexample in Theorem~\ref{thm:unbounded} was solved with
original optimal value $-2$ and a projected unbounded status.
For the forward implication, we used
\[
A=\begin{pmatrix}1&-1&0\\0&0&1\end{pmatrix},\qquad
b=\begin{pmatrix}1\\1\end{pmatrix},\qquad
c=\begin{pmatrix}-1\\0\\0\end{pmatrix},\qquad
S=\begin{pmatrix}1&1\end{pmatrix}.
\]
Here $x_0=(1,0,1)^\top$ and $d=(1,1,0)^\top$ satisfy
$Ax_0=b$, $Ad=0$, $d\geq0$, and $c^\top d=-1$.
Both original and projected LPs were reported unbounded.
These checks illustrate the theorem; its proof establishes the
implication for all admissible LPs.

An additional full-dimensional invertible projection returned the same
optimal value as its original LP. We also checked the equality-dual sign
convention, reduced-cost nonnegativity, and complementary slackness on
that deterministic problem. These checks guard against implementation
errors in the projected solve and lifted-dual recovery.

The new repository contains the independent implementation, pinned
numerical dependencies, raw and aggregate results, source checksums,
environment metadata, deterministic checks, and separate solver
diagnostics. A verifier recomputes the aggregate statistics from the raw
records and checks pairing, counts, provenance hashes, and diagnostic
coverage. The report tables are generated from those aggregates and can
be checked against the committed results. A rerun uses a new output
directory to retain the recorded baseline. Timing values can change
with hardware and system load.

\section{Summary of the New Experiments}
% BEGIN GENERATED key_results
The fresh run comprises 300 projected LPs from 180 independent originals. Dual and primal recoveries yielded 0/60 and 0/60 accepted original-feasible candidates, respectively. Feasible-grid time ratios range from $1.07$ to $4.07$, measuring relaxation computation rather than successful original-LP solution. All-negative projections preserved infeasibility in 120/120 cases. The mixed-sign stress test produced 99 infeasible, 17 optimal and 4 unresolved error outcomes.
% END GENERATED key_results

\chapter{Netlib Benchmark Extension}
\section{Scope and Reproducible Protocol}
We extend the synthetic evaluation using five established LPs from the
Netlib collection \cite{netlib}. BLEND is an oil-refinery model variant;
STOCFOR1 is a deterministic member of a forestry-planning family. AFIRO
and ADLITTLE were supplied by Stanford SOL, and SCFXM1 belongs to the
Ho--Loute staircase collection. These are historical application-derived
or established solver benchmarks, not newly collected field data.
The selection is a small, fixed benchmark suite, not a representative
sample of all industrial LPs.

The source files are downloaded from a pinned COIN-OR mirror revision
and checked by SHA-256 \cite{netlibmirror}. Each original row inequality
is converted exactly to an equality with a nonnegative slack or surplus
variable. The selected models have nonnegative variables with no finite
upper bounds. Minimization sense and zero objective offset are checked;
unsupported model types are rejected. Both the native and standard-form
LPs match Netlib's published optima to the stated numerical tolerance.
We do not add artificial variable bounds to force the projected LPs to
remain bounded.

% BEGIN GENERATED netlib_datasets
\par\medskip
\noindent\begin{minipage}{\linewidth}
\small\linespread{1}\selectfont\centering
\setlength{\tabcolsep}{4pt}
\renewcommand{\arraystretch}{1.08}
\captionof{table}{Benchmark dimensions after exact standard-form conversion. Original storage is CSR; objective values are those recomputed by HiGHS.}\label{tab:netlib-datasets}
\medskip
\sbox0{%
\begin{tabular}{l rrr rr}
\toprule
LP & $m$ & $n$ & $\operatorname{nnz}(A)$ & CSR (KiB) & Optimum \\
\midrule
ADLITTLE & 56 & 138 & 424 & 5.19 & 225494.963162 \\
AFIRO & 27 & 51 & 102 & 1.30 & -464.753143 \\
BLEND & 74 & 114 & 522 & 6.41 & -30.812150 \\
SCFXM1 & 330 & 600 & 2732 & 33.31 & 18416.759028 \\
STOCFOR1 & 117 & 165 & 501 & 6.33 & -41131.976219 \\
\bottomrule
\end{tabular}
}%
\ifdim\wd0>\linewidth
  \resizebox{\linewidth}{!}{\usebox0}
\else
  \usebox0
\fi
\end{minipage}
\par\medskip
% END GENERATED netlib_datasets

For each of these five \emph{fixed} LPs we use Gaussian, sparse-valued,
and CountSketch projectors, $k=\lfloor\rho m\rfloor$ for
$\rho\in\{0.5,0.75\}$, and ten independent projector seeds. Thus there
are 300 projected solves, but only five distinct original LPs. Variability
within a configuration is over sketches, not independently sampled
application problems. Projector distributions follow the preceding
definitions. The master seed is \texttt{20261004}; independent streams
are derived from dataset, ratio, method, and repetition indices.
Method execution order rotates across repetitions.

All timed solves use SciPy's HiGHS dual-simplex method, presolve, feasibility
tolerances $10^{-7}$, a single requested thread, and a 30-second per-solve
limit. The run used Python 3.12.14, NumPy 2.5.3, SciPy 1.18.1, and HiGHS
1.12.0 on macOS arm64, on 4 October 2026. An original solve is paired
with each dataset/ratio/repetition and shared by its three projector
comparisons. Projected time includes sampling, multiplication and sparse
format conversion, and the solver call. Recovery is timed separately.
Reading/converting MPS input, quality checks, and certificate solves are
excluded from benchmark timings. Every group has ten observations;
raw times, sample standard deviations, and conditional metric counts
are supplied in the repository.

\section{Solver Status and Original Feasibility}
Because all five original LPs are feasible and bounded, their projected
relaxations are feasible in exact arithmetic, but need not remain bounded.
For each reported unbounded projection we separately solve for a
nonnegative direction $d$ satisfying
\[
SA d=0,\qquad c^\top d=-1.
\]
The stored numerical certificate must have
$\norm{SA d}_\infty\leq10^{-6}$, $\min_i d_i\geq-10^{-7}$, and
$|c^\top d+1|\leq10^{-6}$. Together with an original feasible point,
this gives a numerical witness of projected unboundedness. The
certificate is an untimed diagnostic, not part of the claimed workflow
cost. Its original residual $\norm{Ad}_\infty$ is also recorded.

% BEGIN GENERATED netlib_status
\par\medskip
\noindent\begin{minipage}{\linewidth}
\small\linespread{1}\selectfont\centering
\setlength{\tabcolsep}{4pt}
\renewcommand{\arraystretch}{1.08}
\captionof{table}{Projected status counts, ten sketches per cell pair. O: finite optimum; U: unbounded. No infeasible, limit, or error outcomes occurred.}\label{tab:netlib-status}
\medskip
\sbox0{%
\begin{tabular}{ll rrrr}
\toprule
LP & Projector & $\rho=.5$: O & U & $\rho=.75$: O & U \\
\midrule
ADLITTLE & Gaussian & 4 & 6 & 10 & 0 \\
ADLITTLE & Sparse & 3 & 7 & 10 & 0 \\
ADLITTLE & CountSketch & 0 & 10 & 1 & 9 \\
AFIRO & Gaussian & 1 & 9 & 9 & 1 \\
AFIRO & Sparse & 0 & 10 & 8 & 2 \\
AFIRO & CountSketch & 0 & 10 & 1 & 9 \\
BLEND & Gaussian & 0 & 10 & 8 & 2 \\
BLEND & Sparse & 0 & 10 & 9 & 1 \\
BLEND & CountSketch & 0 & 10 & 0 & 10 \\
SCFXM1 & Gaussian & 0 & 10 & 10 & 0 \\
SCFXM1 & Sparse & 0 & 10 & 10 & 0 \\
SCFXM1 & CountSketch & 0 & 10 & 0 & 10 \\
STOCFOR1 & Gaussian & 0 & 10 & 0 & 10 \\
STOCFOR1 & Sparse & 0 & 10 & 0 & 10 \\
STOCFOR1 & CountSketch & 0 & 10 & 0 & 10 \\
\bottomrule
\end{tabular}
}%
\ifdim\wd0>\linewidth
  \resizebox{\linewidth}{!}{\usebox0}
\else
  \usebox0
\fi
\end{minipage}
\par\medskip
% END GENERATED netlib_status

% BEGIN GENERATED netlib_findings
Across 300 projected LPs, 216 are unbounded and 84 have finite optima. All 216 unbounded outcomes have passing numerical ray certificates; none of the stored directions satisfies the original equalities at the certificate tolerance. Of the 84 finite projected optima, 0 raw vectors pass original feasibility. Dual recovery encounters a singular selected basis in 84/84 attempts and produces no candidate. Primal recovery produces 84 candidates, all from rank-deficient selected systems (84/84); 0 pass original feasibility. These are failures of the specific recovery heuristics, not evidence that no feasible retrieval method exists.
% END GENERATED netlib_findings

The acceptance rule is the same as in the synthetic experiments: both
the original equality residual and nonnegativity must pass. For raw
projected vectors, the native MPS inequalities are checked as well;
none passes. The two thesis recovery procedures are applied only when
the projected LP has a finite optimum. A singular selected basis is
reported as a recovery failure, not silently dropped from its attempts.
An invertible, full-dimensional orthogonal control ($k=m$) is also run
on each LP: all five preserve the original optimum and feasibility.
This isolates the effect of constraint reduction from a mere change of
coordinates.

\section{Time and Storage Trade-offs}
Unlike the synthetic suite, the Netlib originals are stored as CSR
matrices. Gaussian products are dense; sparse-valued and CountSketch
products are stored in CSR. ``Matrix storage'' below counts the actual
array buffers (values, indices, and row pointers for CSR), not peak
process memory. Projector buffers are recorded separately in the raw
data. The legacy recovery heuristics require a dense copy of $A$;
its storage is explicitly excluded from the projected-matrix comparison,
so these figures are not total end-to-end memory savings.

% BEGIN GENERATED netlib_resources
\par\medskip
\noindent\begin{minipage}{\linewidth}
\small\linespread{1}\selectfont\centering
\setlength{\tabcolsep}{4pt}
\renewcommand{\arraystretch}{1.08}
\captionof{table}{Mean time and projected-matrix storage at $\rho=0.75$ (ten sketches per row). Original/projected times are milliseconds; projected totals include sampling and multiplication. Density is the percentage of nonzero matrix entries. All statuses are included.}\label{tab:netlib-resources}
\medskip
\sbox0{%
\begin{tabular}{ll rrrr}
\toprule
LP & Projector & Orig. (ms) & Proj. (ms) & Density (\%) & $SA$ (KiB) \\
\midrule
ADLITTLE & Gaussian & 1.41 & 3.53 & 100.0 & 45.28 \\
ADLITTLE & Sparse & 1.41 & 3.55 & 61.0 & 41.58 \\
ADLITTLE & CountSketch & 1.41 & 1.21 & 7.0 & 4.92 \\
AFIRO & Gaussian & 0.77 & 0.91 & 100.0 & 7.97 \\
AFIRO & Sparse & 0.77 & 1.01 & 50.5 & 6.12 \\
AFIRO & CountSketch & 0.77 & 0.77 & 9.5 & 1.22 \\
BLEND & Gaussian & 1.50 & 4.58 & 100.0 & 48.98 \\
BLEND & Sparse & 1.50 & 4.01 & 64.8 & 47.82 \\
BLEND & CountSketch & 1.50 & 1.32 & 7.8 & 5.94 \\
SCFXM1 & Gaussian & 4.89 & 968.90 & 100.0 & 1157.81 \\
SCFXM1 & Sparse & 4.89 & 481.99 & 63.6 & 1105.41 \\
SCFXM1 & CountSketch & 4.89 & 5.86 & 1.8 & 32.51 \\
STOCFOR1 & Gaussian & 1.40 & 13.23 & 100.0 & 112.15 \\
STOCFOR1 & Sparse & 1.40 & 11.62 & 60.5 & 102.04 \\
STOCFOR1 & CountSketch & 1.40 & 1.52 & 3.4 & 6.12 \\
\bottomrule
\end{tabular}
}%
\ifdim\wd0>\linewidth
  \resizebox{\linewidth}{!}{\usebox0}
\else
  \usebox0
\fi
\end{minipage}
\par\medskip
% END GENERATED netlib_resources

The table reports means at $\rho=0.75$; both ratios and all dispersion
measures remain available in machine-readable form. Times include all
statuses, including unbounded outcomes. A smaller or faster projected
problem is not a faster solution of the original LP when it is
unbounded or its recovered vector is infeasible. CountSketch can
retain much smaller matrix buffers, but in this experiment that storage
benefit does not translate into an accepted original-feasible solution.

\section{What Is Extended Relative to the Thesis?}
Computational limitations were already discussed in the thesis, not
discovered for the first time here. Its experimental discussion
\cite[pp.~38--40]{thesis2025} identifies projection overhead, densification
of $SA$, dense QR costs, and the efficiency of the original solver's
presolve. Table 4.1 also reports substantial negative mass even where
equality residuals round to zero. The interior-point discussion on
pp.~31--35 concerns time and memory costs in the literature, not a
measured hardware limit in the author's experiments. We found no explicit
out-of-memory event or stated RAM cap explaining the largest tested size.

The new contribution to this course report is therefore an empirical and
reproducibility extension: explicit acceptance counts, status separation,
paired seeded experiments, application-related benchmarks, real sparse
storage, CountSketch, two compression ratios, saved solution vectors,
and checkable unbounded-ray certificates. It does not establish an
asymptotic improvement or a new recovery algorithm. The original thesis
prose about projected feasibility is stronger than its all-zero
misclassification tables support; the new mixed-sign family tests that
possibility separately, without rewriting the old evidence.

The benchmarks are deliberately modest in size and run in one software
environment. They do not test large-scale memory exhaustion, provide a
scaling law, or exhaust retrieval methods. Improved retrieval procedures
exist, including the work of Liberti, Manca, and Poirion
\cite{liberti2023}; they are not evaluated here. All conclusions are
restricted to the basic relaxation and the two stated thesis heuristics.

The repository supplies raw results, saved vectors, version pins and hashes.
Its read-only verifier reconstructs sketches, checks metrics, certificates
and summaries, and tests row conversion and determinism. Generated tables
are checked against the results. Raw MPS files are fetched, not redistributed:
the mirror's build-system license excludes them. The synthetic baseline
is preserved unchanged.

\chapter{Conclusion}
This project explored matrix sketching for linear programming, using
the author's master's thesis as its theoretical starting point.
The Johnson--Lindenstrauss lemma motivates dimensionality reduction
through distance preservation for finite sets. The retained
unboundedness result follows from a different property: projecting
equalities enlarges the feasible set. An original unbounded ray remains
feasible, whereas the converse implication fails.

The revised synthetic evaluation did not yield valid recovered solutions
under the two tested heuristics. All legacy all-negative projections
remained infeasible, whereas the additional mixed-sign family produced
some feasible relaxations and unresolved primary statuses. The instance
distribution and recovery procedure therefore matter alongside running time.

The Netlib extension adds five fixed benchmarks, two compression ratios,
and three projectors, including CountSketch. Of 300 projections, 216
are unbounded with verified numerical ray certificates, while 84 have
finite optima; neither the raw nor the recovered candidates pass original
feasibility. All five full-dimensional controls pass. The covariance
and subspace-embedding discussion clarifies why geometric sketch quality
alone does not imply validity of a compressed equality LP.

These small benchmarks use one solver environment and do not establish
large-scale performance guarantees. Smaller matrix storage or faster
relaxation computation is not a successful LP speedup without an accepted
original-feasible solution. Other retrieval and sketch-based optimization
methods remain outside this evaluation.

\begin{thebibliography}{99}
\bibitem{vu2018}
K. Vu, P.-L. Poirion, and L. Liberti.
Random projections for linear programming.
\emph{Mathematics of Operations Research}, 43(4):1051--1071, 2018.
\href{https://doi.org/10.1287/moor.2017.0894}{doi:10.1287/moor.2017.0894}.

\bibitem{dasgupta2003}
S. Dasgupta and A. Gupta.
An elementary proof of a theorem of Johnson and Lindenstrauss.
\emph{Random Structures \& Algorithms}, 22(1):60--65, 2003.
\href{https://doi.org/10.1002/rsa.10073}{doi:10.1002/rsa.10073}.

\bibitem{achlioptas2003}
D. Achlioptas.
Database-friendly random projections: Johnson--Lindenstrauss with binary coins.
\emph{Journal of Computer and System Sciences}, 66(4):671--687, 2003.

\bibitem{vanderbei2020}
R. J. Vanderbei.
\emph{Linear Programming: Foundations and Extensions}.
5th edition, Springer, 2020.

\bibitem{thesis2025}
J. Zhumabek.
\emph{Matrix Sketching and Linear Programming}.
M.Sc. thesis in Applied Mathematics, Nazarbayev University, 2025.
\bibitem{legacyrepo}
J. Zhumabek.
\emph{Random\_Projections\_master\_thesis}.
Companion code for the master's thesis. GitHub repository, 2025.
\url{https://github.com/Jamil997/Random_Projections_master_thesis}.
Inspected at commit \texttt{3d2ff96}; rechecked 4 October 2026.

\bibitem{clarkson2013}
K. L. Clarkson and D. P. Woodruff.
Low rank approximation and regression in input sparsity time.
\emph{Proceedings of STOC}, pp.~81--90, 2013.
\url{https://arxiv.org/abs/1207.6365}.

\bibitem{halko2011}
N. Halko, P.-G. Martinsson, and J. A. Tropp.
Finding structure with randomness: Probabilistic algorithms for
constructing approximate matrix decompositions.
\emph{SIAM Review}, 53(2):217--288, 2011.
\url{https://arxiv.org/abs/0909.4061}.

\bibitem{ghashami2016}
M. Ghashami, E. Liberty, J. M. Phillips, and D. P. Woodruff.
Frequent Directions: Simple and deterministic matrix sketching.
\emph{SIAM Journal on Computing}, 45(5):1762--1792, 2016.
\url{https://arxiv.org/abs/1501.01711}.

\bibitem{chowdhury2022}
A. Chowdhury, G. Dexter, P. London, H. Avron, and P. Drineas.
Faster randomized interior point methods for tall/wide linear programs.
\emph{Journal of Machine Learning Research}, 23(336):1--48, 2022.
\url{https://jmlr.org/papers/v23/21-0709.html}.

\bibitem{liberti2023}
L. Liberti, B. Manca, and P.-L. Poirion.
Random projections for linear programming: An improved retrieval phase.
\emph{ACM Journal of Experimental Algorithmics}, 28, Article 2.2, 2023.
\href{https://doi.org/10.1145/3617506}{doi:10.1145/3617506}.

\bibitem{netlib}
Netlib. Linear programming test problems: descriptions and reference optima.
\url{https://www.netlib.org/lp/data/readme}. Accessed 4 October 2026.

\bibitem{netlibmirror}
COIN-OR. \emph{Data-Netlib}, MPS benchmark mirror.
\url{https://github.com/coin-or-tools/Data-Netlib}.
Revision \texttt{f1cc423067407d55d579c9c35fb01edf860dbc24}.
\end{thebibliography}
\end{document}
